\documentclass[12pt,a4paper]{article}

\usepackage[a4paper,margin=1in,headheight=15pt]{geometry}
\usepackage{graphicx}
\usepackage{booktabs}
\usepackage{amsmath,amssymb}
\usepackage{float}
\usepackage{hyperref}
\usepackage{xcolor}
\usepackage{listings}
\usepackage{enumitem}
\usepackage{fancyhdr}
\usepackage{longtable}
\usepackage{array}
\usepackage{multirow}
\usepackage{subcaption}
\usepackage{caption}
\usepackage{url}
\Urlmuskip=0mu plus 2mu\relax
\sloppy

\hypersetup{colorlinks=true,linkcolor=blue,urlcolor=blue}

\pagestyle{fancy}
\fancyhf{}
\lhead{ICS1512}
\rhead{Machine Learning Laboratory}
\cfoot{\thepage}

\definecolor{codebg}{RGB}{245,245,245}
\definecolor{keywordcolor}{RGB}{0,0,180}
\definecolor{commentcolor}{RGB}{80,120,80}
\definecolor{stringcolor}{RGB}{163,21,21}

\lstset{
  language=Python,
  basicstyle=\ttfamily\footnotesize,
  numbers=left,
  numberstyle=\tiny\color{gray},
  frame=single,
  breaklines=true,
  showstringspaces=false,
  backgroundcolor=\color{codebg},
  keywordstyle=\color{keywordcolor}\bfseries,
  commentstyle=\color{commentcolor}\itshape,
  stringstyle=\color{stringcolor},
  tabsize=4,
  captionpos=b
}

\begin{document}

% ─────────────────────────────────────────────────────────────────────────────
%  COVER TABLE
% ─────────────────────────────────────────────────────────────────────────────
\begin{center}
  \textbf{\large Sri Sivasubramaniya Nadar College of Engineering, Chennai}\\[4pt]
  \textit{An Autonomous Institution Affiliated to Anna University}\\[12pt]
\end{center}

\begin{tabular}{|p{4cm}|p{11cm}|}
\hline
Experiment No. & 6 \\ \hline
Experiment Title & Bagging, Boosting, and Stacked Ensemble Models \\ \hline
Student Name & Mohamed Farih \\ \hline
Register Number & 3122247001035 \\ \hline
GitHub Repository & \url{https://github.com/farih07/Machine-Learning-Lab} \\ \hline
Faculty & Dr. Poreddy Ajay Kumar Reddy \\ \hline
Submission Date & September 21,2026 \\ \hline
\end{tabular}

\vspace{1cm}

% ─────────────────────────────────────────────────────────────────────────────
\section{Objective}
% ─────────────────────────────────────────────────────────────────────────────
\begin{itemize}[noitemsep]
  \item To understand ensemble learning strategies: Bagging, Boosting, and Stacking.
  \item To implement Bagging and Boosting (AdaBoost, Gradient Boosting) classifiers.
  \item To build a Stacked Ensemble model using multiple heterogeneous base learners.
  \item To compare ensemble models in terms of accuracy, stability, and generalisation.
  \item To analyse the effect of ensemble methods on bias and variance.
\end{itemize}

% ─────────────────────────────────────────────────────────────────────────────
\section{Problem Statement}
% ─────────────────────────────────────────────────────────────────────────────
Given 30 real-valued features derived from digitised fine-needle aspirate (FNA) images of breast
masses, the objective is to classify each sample as \textbf{Malignant} (class 0) or \textbf{Benign}
(class 1) using three distinct ensemble strategies. The input feature vector
$\mathbf{x} \in \mathbb{R}^{30}$ is processed by Bagging, Boosting, and Stacking classifiers whose
hyperparameters are tuned via 5-fold stratified cross-validation. Models are evaluated on a held-out
20\% test split to determine which ensemble strategy achieves the best diagnostic performance while
preserving generalisability.

% ─────────────────────────────────────────────────────────────────────────────
\section{Dataset Description}
% ─────────────────────────────────────────────────────────────────────────────

\begin{longtable}{|p{5.5cm}|p{9cm}|}
\hline
\textbf{Attribute}      & \textbf{Detail} \\ \hline
Dataset Name            & Wisconsin Diagnostic Breast Cancer (WDBC) \\ \hline
Dataset Source          & UCI Machine Learning Repository (\url{https://archive.ics.uci.edu}) \\ \hline
Number of Samples       & 569 \\ \hline
Number of Features      & 30 numerical attributes (mean, SE, worst of 10 nucleus measurements) \\ \hline
Number of Classes       & 2 --- Malignant (M): 212, Benign (B): 357 \\ \hline
Missing Values          & None \\ \hline
Train--Test Split       & 80\% training (455 samples) / 20\% test (114 samples), stratified \\ \hline
Label Encoding          & Malignant $\rightarrow$ 0, Benign $\rightarrow$ 1 \\ \hline
\end{longtable}

% ─────────────────────────────────────────────────────────────────────────────
\section{Exploratory Data Analysis}
% ─────────────────────────────────────────────────────────────────────────────

\subsection*{4.1 Class Distribution}

\begin{figure}[H]
\centering
\includegraphics[width=0.52\linewidth]{exp6_figures/e7_fig1_class}
\caption{Class distribution of the WDBC dataset.
  Benign samples (357) outnumber malignant samples (212) at a ratio of approximately 1.68:1.}
\end{figure}

\textbf{Inference:}
The dataset exhibits a moderate class imbalance (Benign:Malignant $\approx$ 1.68:1), which
motivates the use of stratified sampling throughout training and cross-validation. Since no
resampling technique is applied, models are implicitly evaluated on a realistic clinical distribution.
The minority malignant class is clinically critical --- false negatives (missed malignancies) carry
a higher cost than false positives, making recall for class 0 a key diagnostic metric.

\subsection*{4.2 Feature Correlation Analysis}

\begin{figure}[H]
\centering
\includegraphics[width=0.88\linewidth]{exp6_figures/e7_fig2_corr}
\caption{Pairwise Pearson correlation heatmap for the 12 highest-variance features.
  Saturated red/blue denotes strong positive/negative correlation.}
\end{figure}

\textbf{Inference:}
Strong positive correlations (|r| $>$ 0.9) exist among the \textit{radius}, \textit{perimeter},
and \textit{area} feature groups, which are geometrically linked ($A \approx \pi r^2$).
This multicollinearity has limited effect on tree-based ensembles (which split greedily on
individual features), but it benefits the Stacking model by allowing diverse base learners
(SVM, Decision Tree) to exploit different aspects of the feature space. Concavity-related
features carry the strongest discriminative signal between malignant and benign classes.

% ─────────────────────────────────────────────────────────────────────────────
\section{Data Preprocessing}
% ─────────────────────────────────────────────────────────────────────────────

\begin{itemize}[noitemsep]
  \item \textbf{Missing values:} None present; no imputation required.
  \item \textbf{Label encoding:} Categorical labels M/B are encoded as integers
        (M$\rightarrow$0, B$\rightarrow$1) by \texttt{sklearn}'s built-in loader.
  \item \textbf{Feature scaling:} Applied only inside the Stacking pipeline for the SVM base
        learner via \texttt{StandardScaler} wrapped in a \texttt{Pipeline}.
        Bagging, AdaBoost, and Gradient Boosting are scale-invariant.
  \item \textbf{Train--test split:} 80--20 stratified split with \texttt{random\_state=42}.
  \item \textbf{Feature engineering:} No derived features; all 30 raw attributes used directly.
\end{itemize}

% ─────────────────────────────────────────────────────────────────────────────
\section{Machine Learning Algorithms}
% ─────────────────────────────────────────────────────────────────────────────

\subsection*{6.1 Bagging (Bootstrap Aggregation)}

\textbf{Working Principle:}
Bagging trains $T$ independent base classifiers, each on a bootstrapped resample of the training
data (drawn with replacement). Predictions are aggregated by majority vote.

\textbf{Mathematical Formulation:}
Each bootstrap sample $\mathcal{D}_t \sim \mathcal{D}$ with $|\mathcal{D}_t| = N$.
The ensemble prediction is:
\[
\hat{y}_{\text{bag}} = \underset{k}{\arg\max} \sum_{t=1}^{T} \mathbf{1}[h_t(\mathbf{x}) = k]
\]
Bagging reduces generalisation error primarily by reducing \textbf{variance}:
\[
\text{Var}[\hat{y}_{\text{bag}}] \approx \rho \sigma^2 + \frac{1-\rho}{T} \sigma^2
\]
where $\rho$ is the inter-tree correlation and $\sigma^2$ is the single-tree variance.

\textbf{Advantages:} Robust to overfitting; parallelisable; stable predictions.\\
\textbf{Limitations:} Limited improvement on low-variance models; does not reduce bias.

\subsection*{6.2 Boosting}

\textbf{Working Principle:}
Boosting trains models sequentially, each focusing on the misclassified samples of the previous model.

\paragraph{AdaBoost:}
Each weak learner $h_t$ is trained on a reweighted distribution $\mathcal{D}_t$.
After training, the learner receives weight $\alpha_t$ inversely proportional to its error $\epsilon_t$:
\[
\alpha_t = \frac{1}{2} \ln\!\left(\frac{1 - \epsilon_t}{\epsilon_t}\right)
\]
Sample weights are updated: $w_{t+1}(i) \propto w_t(i) \exp(-\alpha_t y_i h_t(\mathbf{x}_i))$.
The final prediction is:
\[
H(\mathbf{x}) = \text{sign}\!\left(\sum_{t=1}^{T} \alpha_t h_t(\mathbf{x})\right)
\]

\paragraph{Gradient Boosting:}
Each new tree $h_t$ is fit to the negative gradient of the loss function with respect to the
current ensemble prediction. For log-loss:
\[
F_t(\mathbf{x}) = F_{t-1}(\mathbf{x}) + \eta \cdot h_t(\mathbf{x})
\]
where $\eta$ is the learning rate. This reduces \textbf{bias} by iteratively correcting residuals.

\textbf{Advantages:} High accuracy; flexible loss functions; handles heterogeneous data.\\
\textbf{Limitations:} Prone to overfitting; sensitive to noisy data; sequential training is slow.

\subsection*{6.3 Stacked Ensemble (Stacking)}

\textbf{Working Principle:}
Stacking trains multiple heterogeneous base learners on the training data, then trains a meta-learner
on their out-of-fold predictions to learn an optimal combination.

\textbf{Mathematical Formulation:}
Let $\{h_1, h_2, \ldots, h_K\}$ be base learners. For each sample $\mathbf{x}_i$, a meta-feature
vector is constructed:
\[
\tilde{\mathbf{x}}_i = [h_1(\mathbf{x}_i),\; h_2(\mathbf{x}_i),\; \ldots,\; h_K(\mathbf{x}_i)]
\]
The meta-learner $g$ (Logistic Regression here) is trained on $\tilde{\mathbf{X}}$:
\[
\hat{y} = g(\tilde{\mathbf{x}}) = \sigma\!\left(\mathbf{w}^\top \tilde{\mathbf{x}} + b\right)
\]

\textbf{Base Models:} SVM (with StandardScaler), Decision Tree ($\texttt{max\_depth}=5$).\\
\textbf{Meta-Learner:} Logistic Regression (\texttt{max\_iter=1000}).

\textbf{Advantages:} Exploits complementary strengths of diverse models; often achieves top accuracy.\\
\textbf{Limitations:} Computationally expensive; risk of data leakage if not handled with CV.

% ─────────────────────────────────────────────────────────────────────────────
\section{Experimental Setup}
% ─────────────────────────────────────────────────────────────────────────────

\begin{tabular}{ll}
\toprule
\textbf{Component}        & \textbf{Detail} \\
\midrule
Python Version            & 3.12 \\
Scikit-learn Version      & 1.8.0 \\
NumPy / Pandas            & 1.26.x / 2.2.x \\
Matplotlib / Seaborn      & 3.9.x / 0.13.x \\
IDE                       & Jupyter Notebook / VS Code \\
Operating System          & Ubuntu 22.04 LTS \\
Hardware                  & CPU: Intel Core i5 (8-core), RAM: 16 GB \\
\bottomrule
\end{tabular}

% ─────────────────────────────────────────────────────────────────────────────
\section{Hyperparameter Selection via 5-Fold Cross-Validation}
% ─────────────────────────────────────────────────────────────────────────────

All hyperparameters were selected by exhaustive grid search evaluated under
\texttt{StratifiedKFold(n\_splits=5, shuffle=True, random\_state=42)}.
The configuration maximising mean CV accuracy was retained for final evaluation.

\subsection*{8.1 Bagging — Search Space}

\begin{center}
\begin{tabular}{ll}
\toprule
\textbf{Hyperparameter} & \textbf{Candidates} \\
\midrule
\texttt{n\_estimators}   & 25, 50, 100, 200 \\
\texttt{max\_samples}    & 0.7, 0.9, 1.0 \\
\texttt{max\_features}   & 0.7, 1.0 \\
Base estimator           & \texttt{DecisionTreeClassifier} \\
\bottomrule
\end{tabular}
\end{center}

\subsection*{Table 1 — Bagging Hyperparameter Evaluation (Top 8 by CV Accuracy)}

\begin{center}
\begin{tabular}{ccccc}
\toprule
\textbf{n\_estimators} & \textbf{max\_samples} & \textbf{max\_features}
  & \textbf{Avg CV Acc (\%)} & \textbf{Avg CV F1} \\
\midrule
25  & 0.9 & 1.0 & 96.26 & 0.9702 \\
50  & 0.9 & 1.0 & 96.26 & 0.9701 \\
100 & 0.9 & 0.7 & 96.04 & 0.9683 \\
200 & 0.7 & 0.7 & 96.04 & 0.9683 \\
25  & 0.9 & 0.7 & 95.82 & 0.9667 \\
100 & 1.0 & 1.0 & 95.60 & 0.9647 \\
200 & 0.9 & 1.0 & 95.60 & 0.9647 \\
50  & 0.7 & 1.0 & 95.38 & 0.9631 \\
\bottomrule
\end{tabular}
\end{center}

\textbf{Selected Bagging Configuration:}
\texttt{n\_estimators=50, max\_samples=0.9, max\_features=1.0}
(Best avg CV Accuracy: \textbf{96.26\%})

\subsection*{8.2 Boosting — Search Space}

\begin{center}
\begin{tabular}{lll}
\toprule
\textbf{Model}        & \textbf{Hyperparameter}  & \textbf{Candidates} \\
\midrule
\multirow{2}{*}{AdaBoost} & \texttt{n\_estimators} & 50, 100, 200 \\
                           & \texttt{learning\_rate}& 0.1, 0.5, 1.0 \\
\midrule
\multirow{3}{*}{Gradient Boosting} & \texttt{n\_estimators} & 50, 100, 150 \\
                                   & \texttt{learning\_rate}& 0.05, 0.1, 0.2 \\
                                   & \texttt{max\_depth}    & 2, 3 \\
\bottomrule
\end{tabular}
\end{center}

\subsection*{Table 2 — Boosting Hyperparameter Evaluation (Top 8 — AdaBoost \& Gradient Boosting)}

\textbf{AdaBoost:}
\begin{center}
\begin{tabular}{cccc}
\toprule
\textbf{n\_estimators} & \textbf{learning\_rate}
  & \textbf{Avg CV Acc (\%)} & \textbf{Avg CV F1} \\
\midrule
100 & 1.0 & 98.02 & 0.9844 \\
200 & 1.0 & 97.80 & 0.9826 \\
100 & 0.5 & 97.36 & 0.9792 \\
200 & 0.1 & 97.14 & 0.9773 \\
200 & 0.5 & 96.92 & 0.9758 \\
50  & 1.0 & 96.70 & 0.9738 \\
50  & 0.5 & 95.60 & 0.9647 \\
50  & 0.1 & 94.51 & 0.9553 \\
\bottomrule
\end{tabular}
\end{center}

\textbf{Selected AdaBoost Configuration:}
\texttt{n\_estimators=100, learning\_rate=1.0}
(Best avg CV Accuracy: \textbf{98.02\%})

\vspace{8pt}
\textbf{Gradient Boosting:}
\begin{center}
\begin{tabular}{ccccc}
\toprule
\textbf{n\_estimators} & \textbf{learning\_rate} & \textbf{max\_depth}
  & \textbf{Avg CV Acc (\%)} & \textbf{Avg CV F1} \\
\midrule
150 & 0.2 & 2 & 97.36 & 0.9789 \\
100 & 0.2 & 2 & 96.92 & 0.9755 \\
150 & 0.2 & 3 & 96.70 & 0.9738 \\
100 & 0.2 & 3 & 96.48 & 0.9719 \\
150 & 0.1 & 2 & 96.48 & 0.9719 \\
100 & 0.1 & 2 & 96.26 & 0.9702 \\
50  & 0.2 & 2 & 96.26 & 0.9702 \\
150 & 0.1 & 3 & 95.82 & 0.9667 \\
\bottomrule
\end{tabular}
\end{center}

\textbf{Selected Gradient Boosting Configuration:}
\texttt{n\_estimators=150, learning\_rate=0.2, max\_depth=2}
(Best avg CV Accuracy: \textbf{97.36\%})

\subsection*{8.3 Stacked Ensemble — Configurations Evaluated}

\subsection*{Table 3 — Stacked Ensemble Hyperparameter Evaluation}

\begin{center}
\begin{tabular}{lccc}
\toprule
\textbf{Base Models}     & \textbf{Meta-Learner}
  & \textbf{Avg CV Acc (\%)} & \textbf{Avg CV F1} \\
\midrule
SVM + DT        & Logistic Regression & 97.14 & 0.9770 \\
SVM + NB + DT   & Logistic Regression & 96.70 & 0.9737 \\
SVM + NB        & Logistic Regression & 96.26 & 0.9701 \\
NB + DT         & Logistic Regression & 93.63 & 0.9500 \\
\bottomrule
\end{tabular}
\end{center}

\textbf{Selected Stacking Configuration:}
Base: SVM (+ StandardScaler) + Decision Tree; Meta-Learner: Logistic Regression
(Best avg CV Accuracy: \textbf{97.14\%})

% ─────────────────────────────────────────────────────────────────────────────
\section{5-Fold Cross-Validation Fold-by-Fold Comparison}
% ─────────────────────────────────────────────────────────────────────────────

\begin{center}
\begin{tabular}{lcccccc}
\toprule
\textbf{Model} & \textbf{Fold 1} & \textbf{Fold 2} & \textbf{Fold 3}
               & \textbf{Fold 4} & \textbf{Fold 5} & \textbf{Average} \\
               & (\%) & (\%) & (\%) & (\%) & (\%) & (\%) \\
\midrule
Bagging           & 96.70 & 95.60 & 94.51 & 95.60 & 98.90 & 96.26 \\
AdaBoost          & 98.90 & 100.00 & 94.51 & 97.80 & 98.90 & 98.02 \\
Gradient Boosting & 95.60 & 100.00 & 94.51 & 97.80 & 98.90 & 97.36 \\
Stacking          & 96.70 & 96.70  & 95.60 & 97.80 & 98.90 & 97.14 \\
\bottomrule
\end{tabular}
\end{center}

\begin{figure}[H]
\centering
\includegraphics[width=0.88\linewidth]{exp6_figures/e7_fig3_cv}
\caption{Per-fold cross-validation accuracy for all four ensemble models.
  AdaBoost achieves the highest mean CV accuracy (98.02\%), but shows higher fold-to-fold variability
  (range: 5.49 pp) compared to Stacking (range: 3.30 pp).}
\end{figure}

\textbf{Inference:}
AdaBoost achieves the best mean CV accuracy (98.02\%) and achieves perfect accuracy on Fold 2
(100\%), demonstrating the power of sequential bias correction. However, it also records the
greatest variability across folds (range: 5.49 pp vs. Stacking's 3.30 pp), suggesting that its
sequential reweighting can overfit to fold-specific noise. Gradient Boosting balances high average
accuracy (97.36\%) with more consistent fold behaviour. Stacking (97.14\%) is the most stable,
confirming that heterogeneous base learners provide complementary error coverage across folds.

% ─────────────────────────────────────────────────────────────────────────────
\section{Results}
% ─────────────────────────────────────────────────────────────────────────────

\subsection*{Table 4 — Performance Comparison of Ensemble Models (Test Set, 80--20 Split)}

\begin{center}
\begin{tabular}{lcccc}
\toprule
\textbf{Model}        & \textbf{Accuracy (\%)} & \textbf{Precision (\%)}
  & \textbf{Recall (\%)} & \textbf{F1-Score (\%)} \\
\midrule
Bagging           & 94.74 & 95.83 & 95.83 & 95.83 \\
AdaBoost          & 95.61 & 94.67 & 98.61 & 96.60 \\
Gradient Boosting & 95.61 & 94.67 & 98.61 & 96.60 \\
\textbf{Stacking} & \textbf{96.49} & \textbf{97.22} & \textbf{97.22} & \textbf{97.22} \\
\bottomrule
\end{tabular}
\end{center}

\begin{center}
\begin{tabular}{lc}
\toprule
\textbf{Model}        & \textbf{AUC-ROC} \\
\midrule
Bagging           & 0.9926 \\
AdaBoost          & 0.9818 \\
Gradient Boosting & 0.9914 \\
\textbf{Stacking} & \textbf{0.9944} \\
\bottomrule
\end{tabular}
\end{center}

Stacking achieves the highest accuracy (96.49\%), F1-score (97.22\%), and AUC-ROC (0.9944) on the
test set, confirming that heterogeneous base-learner diversity translates to the best held-out
performance. AdaBoost and Gradient Boosting achieve equal test accuracy (95.61\%) with high recall
(98.61\%), while Bagging produces the most balanced Precision--Recall trade-off (both at 95.83\%).

% ─────────────────────────────────────────────────────────────────────────────
\section{Result Visualisations}
% ─────────────────────────────────────────────────────────────────────────────

\subsection*{10.1 Confusion Matrix — Bagging}

\begin{figure}[H]
\centering
\includegraphics[width=0.52\linewidth]{exp6_figures/e7_fig4_bag_cm}
\caption{Confusion matrix for the Bagging classifier (DecisionTree base, 50 estimators,
  max\_samples=0.9) on the 114-sample test set.}
\end{figure}

\textbf{Inference:}
Bagging correctly classifies the majority of both classes with balanced precision and recall
(both 95.83\%). The symmetric confusion matrix confirms that bootstrap aggregation of decision
trees effectively reduces their individual high-variance errors.
However, 6 errors remain: 3 false negatives and 3 false positives, indicating
residual variance that sequential methods address more aggressively.

\subsection*{10.2 Confusion Matrix — AdaBoost}

\begin{figure}[H]
\centering
\includegraphics[width=0.52\linewidth]{exp6_figures/e7_fig5_ada_cm}
\caption{Confusion matrix for the AdaBoost classifier (100 stumps, learning rate=1.0)
  on the test set.}
\end{figure}

\textbf{Inference:}
AdaBoost reduces false negatives (missed malignancies) compared to Bagging by concentrating
weight on hard-to-classify boundary samples through its sequential reweighting mechanism.
A recall of 98.61\% for the benign class is achieved, though precision drops slightly to 94.67\%
due to 3 false positives. In clinical terms, AdaBoost's high recall for malignant detection
makes it well-suited for screening applications where false negatives are costly.

\subsection*{10.3 Confusion Matrix — Gradient Boosting}

\begin{figure}[H]
\centering
\includegraphics[width=0.52\linewidth]{exp6_figures/e7_fig6_gb_cm}
\caption{Confusion matrix for the Gradient Boosting classifier
  (150 trees, lr=0.2, max\_depth=2) on the test set.}
\end{figure}

\textbf{Inference:}
Gradient Boosting produces an identical test accuracy and confusion matrix to AdaBoost (95.61\%),
but achieves this through gradient descent on residuals rather than sample reweighting.
The shallow tree depth (\texttt{max\_depth=2}) acts as a strong regulariser, preventing individual
trees from memorising training noise and maintaining the low bias of the boosting chain
without incurring high variance. The identical confusion matrix to AdaBoost confirms that both
boosting variants converge on similar decision boundaries for this dataset.

\subsection*{10.4 Confusion Matrix — Stacking}

\begin{figure}[H]
\centering
\includegraphics[width=0.52\linewidth]{exp6_figures/e7_fig7_stk_cm}
\caption{Confusion matrix for the Stacking classifier (SVM + DT base; Logistic Regression meta)
  on the test set.}
\end{figure}

\textbf{Inference:}
Stacking achieves the lowest total error (4 misclassifications vs.\ 6 for Bagging, 5 for Boosting)
with equal precision and recall of 97.22\%. The SVM base learner captures large-margin linear
boundaries while the Decision Tree captures non-linear splits; the Logistic Regression meta-learner
learns to weight their complementary predictions optimally. This heterogeneous diversity is the key
mechanism behind Stacking's superior test performance.

\subsection*{10.5 ROC Curve Comparison}

\begin{figure}[H]
\centering
\includegraphics[width=0.68\linewidth]{exp6_figures/e7_fig8_roc}
\caption{ROC curves for all four ensemble classifiers.
  AUC values are shown in the legend.}
\end{figure}

\textbf{Inference:}
Stacking achieves the highest AUC-ROC (0.9944), confirming superior probability calibration
across all classification thresholds. Bagging's AUC of 0.9926 is notably high despite a lower
raw accuracy (94.74\%), indicating well-calibrated probability estimates from averaging decision
tree leaf posteriors. AdaBoost's AUC of 0.9818 is lower than expected given its high CV accuracy,
reflecting slight overconfidence in its probability estimates (AdaBoost is not natively
probability-calibrated). Gradient Boosting (0.9914) sits between Bagging and Stacking,
balancing accuracy and probability calibration.

\subsection*{10.6 Test-Set Metrics Comparison}

\begin{figure}[H]
\centering
\includegraphics[width=0.88\linewidth]{exp6_figures/e7_fig9_metrics}
\caption{Side-by-side comparison of Accuracy, Precision, Recall, and F1-Score
  for all ensemble models on the test set.}
\end{figure}

\textbf{Inference:}
Stacking consistently leads across all four metrics. Bagging's precision equals its recall (95.83\%),
reflecting the symmetric aggregation of identically structured trees. AdaBoost and Gradient Boosting
both achieve 98.61\% recall --- the highest among all models --- at the cost of slightly lower
precision (94.67\%), reflecting boosting's focus on correcting false negatives.
The plot confirms that no single ensemble strategy dominates on every metric simultaneously:
the choice depends on the clinical trade-off between precision and recall.

\subsection*{10.7 Gradient Boosting Feature Importances}

\begin{figure}[H]
\centering
\includegraphics[width=0.88\linewidth]{exp6_figures/e7_fig10_fi}
\caption{Top 15 Gradient Boosting feature importances (mean decrease in impurity).
  Darker green bars indicate features exceeding the 0.07 importance threshold.}
\end{figure}

\textbf{Inference:}
The most discriminative features are \textit{worst concave points} (0.1306),
\textit{worst area} (0.1186), and \textit{worst radius} (0.1071) --- all from the ``worst''
measurement group describing the most morphologically extreme cell nucleus in each image.
This aligns with the clinical understanding that malignant cells exhibit enlarged and
irregularly shaped nuclei detectable in FNA images. The top 10 features account for
approximately 84\% of total predictive weight, confirming significant feature redundancy
across the 30-attribute space.

% ─────────────────────────────────────────────────────────────────────────────


% ─────────────────────────────────────────────────────────────────────────────
\section{Discussion}
% ─────────────────────────────────────────────────────────────────────────────

\textbf{How Bagging reduces variance:}
Bagging trains each tree on a different bootstrapped subsample, so each tree memorises a
different slice of the noise distribution. When predictions are aggregated by majority vote,
individual memorisations cancel out, reducing the effective variance without changing the
expected bias of the base learner. In this experiment, Bagging cut decision tree variance
sufficiently to achieve 94.74\% test accuracy --- significantly above a single tree (90.35\%
in Experiment 5).

\textbf{How Boosting addresses bias:}
AdaBoost and Gradient Boosting use very shallow base learners (decision stumps for AdaBoost,
depth-2 trees for Gradient Boosting), which are individually high-bias. Sequential reweighting
(AdaBoost) or gradient descent (Gradient Boosting) iteratively corrects the bias of the ensemble
by focusing resources on misclassified samples. The result is a strong learner with low bias and
controlled variance, achieving 98.02\% mean CV accuracy (AdaBoost) vs.\ 96.26\% for Bagging.

\textbf{Why Stacking benefits from heterogeneous models:}
An SVM and a Decision Tree make errors on different sample regions: the SVM excels at
large-margin linear boundaries; the Decision Tree captures non-linear, axis-aligned splits.
Their error sets have low overlap, so the Logistic Regression meta-learner can learn to trust
the SVM on its confident predictions and the DT on its confident predictions, producing a lower
combined error rate (96.49\% test accuracy) than either base learner alone.

\textbf{Model selection:}
Stacking is the best overall model on this dataset (AUC = 0.9944, F1 = 97.22\%, accuracy = 96.49\%).
For high-recall clinical screening where missing a malignancy is unacceptable, AdaBoost or
Gradient Boosting (recall = 98.61\%) may be preferred despite slightly lower overall accuracy.

% ─────────────────────────────────────────────────────────────────────────────
\section{Observation Questions}
% ─────────────────────────────────────────────────────────────────────────────

\begin{enumerate}

\item \textbf{How does Bagging reduce variance?}\\
  Bagging creates $T$ independently bootstrapped training sets, trains a separate decision tree
  on each, and aggregates via majority vote. Each tree memorises different noise patterns; majority
  voting cancels these independent errors, reducing $\text{Var}[\hat{y}]$ by a factor of
  approximately $\frac{1+(T-1)\rho}{T}$ compared to a single tree, where $\rho$ is the
  inter-tree correlation. The result is a stable ensemble that generalises better than any
  individual tree.

\item \textbf{How does Boosting address model bias?}\\
  Boosting starts from weak learners (stumps or shallow trees) with high bias and low variance.
  Each subsequent learner is trained to correct the residual errors of the running ensemble ---
  either by reweighting misclassified samples (AdaBoost) or fitting the negative gradient of the
  loss (Gradient Boosting). The cumulative effect is a strong learner whose bias decreases
  with each iteration, at the cost of modest variance increase if not regularised by learning
  rate and depth constraints.

\item \textbf{Why does Stacking benefit from heterogeneous models?}\\
  Heterogeneous base learners partition the error space differently: an SVM exploits
  maximum-margin linear boundaries, a Naive Bayes relies on conditional independence, and a
  Decision Tree captures non-linear axis-aligned splits. Their complementary error sets mean the
  meta-learner receives diverse, orthogonal information from each base model's predictions, allowing
  it to combine them more effectively than any homogeneous ensemble. The SVM + DT combination
  achieved the best balance of diversity and accuracy (97.14\% CV).

\item \textbf{Which ensemble method performed best and why?}\\
  \textbf{Stacking} achieved the best overall test performance (accuracy 96.49\%, F1 97.22\%,
  AUC 0.9944). Unlike Bagging and Boosting, Stacking operates over a heterogeneous model space
  and learns an optimal combination function via a meta-learner. On this dataset, where the
  decision boundary has both linear and non-linear components, the SVM + Decision Tree pair
  effectively covers both regions, and the Logistic Regression meta-learner learns to allocate
  confidence between them, achieving the lowest test error.

\end{enumerate}

% ─────────────────────────────────────────────────────────────────────────────
\section{Conclusion}
% ─────────────────────────────────────────────────────────────────────────────

Three distinct ensemble strategies --- Bagging, Boosting (AdaBoost and Gradient Boosting), and
Stacking --- were implemented on the Wisconsin Diagnostic Breast Cancer dataset and tuned via
5-fold stratified cross-validation. The optimal Bagging classifier
(\texttt{n\_estimators=50, max\_samples=0.9}) achieved 96.26\% mean CV accuracy and 94.74\% test
accuracy. AdaBoost (\texttt{n\_estimators=100, learning\_rate=1.0}) achieved the highest mean CV
accuracy (98.02\%) and recall (98.61\%), while Gradient Boosting (\texttt{n\_estimators=150,
lr=0.2, max\_depth=2}) matched AdaBoost's test accuracy with a better AUC (0.9914 vs.\ 0.9818).
Stacking (SVM + DT base; Logistic Regression meta) achieved the best overall test performance
(accuracy 96.49\%, F1 97.22\%, AUC 0.9944), confirming that heterogeneous base-learner
diversity and a learned combination function are the most powerful ensemble strategy for this
classification task.

% ─────────────────────────────────────────────────────────────────────────────
\section{References}
% ─────────────────────────────────────────────────────────────────────────────

\begin{enumerate}[label={[\arabic*]},leftmargin=*]

\item F. Pedregosa \textit{et al.}, ``Scikit-learn: Machine Learning in Python,''
  \textit{Journal of Machine Learning Research}, vol.\ 12, pp.\ 2825--2830, 2011.
  [Online]. Available: \url{https://scikit-learn.org}

\item L. Breiman, ``Bagging Predictors,'' \textit{Machine Learning}, vol.\ 24, no.\ 2,
  pp.\ 123--140, Aug. 1996.
  doi: \href{https://doi.org/10.1007/BF00058655}{10.1007/BF00058655}

\item Y. Freund and R. E. Schapire, ``A Decision-Theoretic Generalization of On-Line Learning
  and an Application to Boosting,'' \textit{Journal of Computer and System Sciences},
  vol.\ 55, no.\ 1, pp.\ 119--139, Aug. 1997.
  doi: \href{https://doi.org/10.1006/jcss.1997.1504}{10.1006/jcss.1997.1504}

\item J. H. Friedman, ``Greedy Function Approximation: A Gradient Boosting Machine,''
  \textit{Annals of Statistics}, vol.\ 29, no.\ 5, pp.\ 1189--1232, 2001.
  doi: \href{https://doi.org/10.1214/aos/1013203451}{10.1214/aos/1013203451}

\item D. H. Wolpert, ``Stacked Generalization,'' \textit{Neural Networks},
  vol.\ 5, no.\ 2, pp.\ 241--259, 1992.
  doi: \href{https://doi.org/10.1016/S0893-6080(05)80023-1}{10.1016/S0893-6080(05)80023-1}

\item D. Dua and C. Graff, ``UCI Machine Learning Repository,''
  University of California, Irvine, 2019.
  [Online]. Available: \url{https://archive.ics.uci.edu/ml/datasets/Breast+Cancer+Wisconsin+(Diagnostic)}

\item Scikit-learn Documentation, ``Ensemble Methods,'' 2024.
  [Online]. Available: \url{https://scikit-learn.org/stable/modules/ensemble.html}

\end{enumerate}

\end{document}